\documentclass[11pt,a4paper]{article}
\usepackage{notes}
\usepackage{fancyhdr}

\pagestyle{fancy}
\fancyhf{}
\lhead{Geophysics IIIA: Potential Fields \& Geothermics}
\rhead{Week 7 --- Tower Hill Case Study}
\cfoot{\thepage}

\title{Week 7 --- Case Study: Reduction to the Pole\\
at the Tower Hill Volcanic Complex, Victoria}
\author{Geophysics IIIA: Potential Fields \& Geothermics}
\date{\today}

\begin{document}
\maketitle

\section{Purpose of the Demonstration}

This case study applies the Fourier-domain magnetic processing methods developed
in the RTP demo notebook to real aeromagnetic data from the
Tower Hill Volcanic Complex in south-western Victoria, Australia.

The demonstration progresses from:
\begin{itemize}
  \item the geometry of a maar/diatreme and why it produces a magnetic anomaly,
  \item the asymmetric character of total magnetic intensity (TMI) anomalies
        in the Southern Hemisphere,
  \item reduction to the pole (RTP) as a Fourier-domain phase correction,
  \item to interpretation of the before- and after-RTP maps.
\end{itemize}

\noindent\textbf{Key message.}
RTP is a phase filter: it does not change wavelengths or add new information.
It shifts and reshapes anomalies so that their maxima overlie their causative
bodies, making maps easier to interpret geologically.

% ─────────────────────────────────────────────────────────────────────────────
\section{Geological Background}

\subsection{Tower Hill Volcanic Complex}

Tower Hill ($38.3^\circ$\,S, $142.4^\circ$\,E) is a Quaternary maar–diatreme complex
located on the Newer Volcanics Province of south-western Victoria.
It formed approximately 32\,000 years ago when rising basaltic magma
encountered groundwater, generating a series of phreatomagmatic explosions.

The dominant structure is a broad maar crater roughly 3.5\,km in diameter,
surrounded by a tuff ring and containing a central scoria cone.
The subsurface consists of a diatreme --- a funnel-shaped breccia pipe ---
filled with highly fragmented, magnetised basaltic material.

\subsection{Why a maar produces a magnetic anomaly}

Mafic magmas are typically rich in magnetite, which becomes permanently
magnetised below the Curie temperature ($\sim$580\,$^{\circ}$C for magnetite)
as the rock cools.
The diatreme fill and associated intrusive material therefore carry both
induced and remanent magnetisation, producing a measurable TMI anomaly
over the crater and its flanks.

In contrast, the surrounding Quaternary sediments and older Palaeozoic basement
have much lower magnetic susceptibility, providing the density contrast that
makes the volcanic complex detectable from the air.

% ─────────────────────────────────────────────────────────────────────────────
\section{The Magnetic Field at Tower Hill}

The International Geomagnetic Reference Field (IGRF) defines the large-scale
inducing field from Earth's core.
At Tower Hill the IGRF parameters are approximately:
\begin{itemize}
  \item Inclination $I \approx -70^\circ$ (field dips steeply into the Earth,
        as expected in the Southern Hemisphere mid-latitudes)
  \item Declination $D \approx +10.5^\circ$ (field points slightly east of
        geographic north)
  \item Total intensity $F \approx 60\,000$\,nT
\end{itemize}

The magnetic anomaly dataset used here is the \emph{IGRF-corrected} total
magnetic intensity: the long-wavelength core field has been removed, leaving
only anomalies due to crustal magnetisation contrasts.

\subsection{Why inclination matters for anomaly shape}

At the magnetic pole ($I = -90^\circ$ in the Southern Hemisphere), the inducing
field is vertical.
The anomaly of a compact, uniformly magnetised body is then symmetric: a
single positive lobe centred above the body.

Away from the pole, the horizontal component of the inducing field
introduces a dipolar signature: the positive lobe is shifted toward the
equator (northward in the Southern Hemisphere) and a negative lobe appears
on the poleward (southern) side of the body.

At Tower Hill ($I \approx -70^\circ$), this asymmetry is pronounced but not
extreme: the positive anomaly maximum is displaced a few kilometres
north of the maar centre, with a weaker negative feature to the south.
Students should verify this displacement in the raw TMI map.

% ─────────────────────────────────────────────────────────────────────────────
\section{Forward Model for Magnetic Anomalies}

\subsection{Thin-sheet spectral approximation}

The total-field magnetic anomaly of a thin magnetised sheet at depth $z$
with thickness distribution $h(x,y)$ is computed in the Fourier domain as
\begin{equation}
\hat{T}(\mathbf{k})
  = \chi\,|\mathbf{k}|\,e^{-|\mathbf{k}|z}\,\Upsilon^2\,\hat{h}(\mathbf{k}),
\end{equation}
where $\mathbf{k} = (k_x, k_y)$ is the horizontal wavenumber vector,
$\chi$ is magnetic susceptibility, and
\begin{equation}
\Upsilon(\mathbf{k})
  = \sin I
    + i\cos I\,
      \frac{k_x\cos D + k_y\sin D}{|\mathbf{k}|}
\end{equation}
is the direction-cosine factor that encodes the inclination and declination
of the inducing field \citep{Blakely1995}.

The exponential factor $e^{-|\mathbf{k}|z}$ is an upward-continuation kernel:
it attenuates short-wavelength content relative to long wavelengths, with
the attenuation rate set by the source depth $z$.
This is why shallow sources produce narrow, high-amplitude anomalies and deep
sources produce broad, smooth anomalies.

\subsection{The role of the direction-cosine factor}

The factor $\Upsilon^2$ in the forward operator is the source of the anomaly
asymmetry.
At the pole ($I = \pm 90^\circ$), $\Upsilon = \pm 1$ (real and constant),
so $|\Upsilon^2|$ is the same for all directions and the anomaly is radially
symmetric.
Away from the pole, $\Upsilon$ becomes complex and direction-dependent:
its phase varies with the angle of $\mathbf{k}$ relative to north, producing
the characteristic dipolar shape of a TMI anomaly.

Reduction to the pole removes exactly this phase: it applies the inverse of
$\Upsilon^2$, replacing it with a factor that is real and direction-independent.

% ─────────────────────────────────────────────────────────────────────────────
\section{Reduction to the Pole}

\subsection{Filter definition}

The RTP filter in the Fourier domain is
\begin{equation}
H_\text{RTP}(\mathbf{k})
  = \frac{\bar{\Upsilon}^2}
         {\bigl(|\Upsilon|^2 + \varepsilon\bigr)^2},
\end{equation}
where $\bar{\Upsilon}$ denotes the complex conjugate of $\Upsilon$,
and $\varepsilon$ is a small water-level stabilisation constant
(typically 0.05 in the notebook) that prevents division by zero.

Applying $H_\text{RTP}$ to $\hat{T}(\mathbf{k})$ is equivalent to replacing
the actual field geometry with a vertical field, transforming the data into
what would have been observed at the magnetic pole.
No new geological information is added; the operation is a re-expression of
the same data.

\subsection{RTP is a phase filter}

Examining the amplitude and phase spectra separately reveals the key property
of RTP: the amplitude spectrum $|H_\text{RTP}|$ is nearly flat (close to 1.0
for most wavenumbers), while the phase of $H_\text{RTP}$ varies strongly with
both wavenumber magnitude and direction.

The demonstration in the notebook makes this explicit:
\begin{itemize}
  \item The log-amplitude maps of $|\hat{T}|$ and $|\hat{T}_\text{RTP}|$ are
        nearly identical.
  \item The phase map of $\hat{T}$ shows a strong N--S asymmetry driven by
        the inclined field.
  \item The phase correction $\angle H_\text{RTP}$ carries the corresponding
        asymmetric pattern that, when added, restores N--S symmetry.
\end{itemize}

\subsection{Instability at low inclinations}

Near the magnetic equator ($I \approx 0^\circ$), the direction-cosine factor
$\Upsilon$ approaches zero for wavenumber components in the N--S direction.
Even with water-level stabilisation, the filter amplifies these components
strongly, producing ringing artefacts and unreliable phase corrections.

Students exploring the inclination slider in the notebook should observe that
artefacts become visible below about $|I| \approx 15$--$20^\circ$.
At Tower Hill ($I \approx -70^\circ$), the field is steeply inclined and well
away from this instability zone, so RTP is reliable.

% ─────────────────────────────────────────────────────────────────────────────
\section{Dataset and Processing Details}

\subsection{Survey parameters}

The Tower Hill dataset consists of an IGRF-corrected aeromagnetic grid
covering a domain of approximately $10.4\,\text{km} \times 13.2\,\text{km}$
centred near $142.36^\circ$\,E, $38.32^\circ$\,S.
The grid is 241\,$\times$\,241 pixels at a spacing of approximately
55\,m in latitude and 43\,m in longitude, giving a Nyquist wavelength of
roughly 110\,m --- far smaller than the maar diameter and the smallest
anomalies of interest.

\subsection{RTP parameters}

Both the ambient field direction and the assumed magnetisation direction
are set to the IGRF values at the site:
$I = -70^\circ$, $D = +10.5^\circ$.
This is the standard assumption of purely induced magnetisation, which is
a reasonable first approximation for Quaternary basalts.
Where remanent magnetisation is significant and differs in direction
from the induced component, the standard RTP will over- or under-correct
individual anomalies; this is discussed further in Section~\ref{sec:limitations}.

\subsection{Edge effects and padding}

The data grid is finite.
Where the FFT treats it as periodic, discontinuities at the grid boundary
can introduce spectral leakage and spatial ringing in the RTP result.
NaN values at the margins of the survey are filled with the grid mean prior
to transformation and restored afterward.
For a real processing workflow, more sophisticated padding (mirroring,
tapering) would be applied.

% ─────────────────────────────────────────────────────────────────────────────
\section{Interpreting the Results}

\subsection{Raw TMI map}

The raw anomaly shows:
\begin{itemize}
  \item a positive lobe displaced slightly north of the maar centre, consistent
        with the southward tilt of the inducing field at $I = -70^\circ$,
  \item a weaker negative lobe to the south of the positive peak,
  \item broadly circular symmetry reflecting the axisymmetric maar/diatreme
        geometry, modified by the field inclination.
\end{itemize}
The anomaly peak amplitude is of order 600--1000\,nT, consistent with a
strongly magnetised basaltic body.

\subsection{RTP map}

After RTP:
\begin{itemize}
  \item the positive lobe is centred approximately over the maar crater,
  \item the negative southern lobe is reduced or eliminated for a predominantly
        induced source,
  \item the anomaly shape is more nearly circular, reflecting the underlying
        axisymmetric geometry of the diatreme.
\end{itemize}

The N--S cross-section through the anomaly maximum (shown in the notebook)
illustrates the northward shift clearly: the peak moves to lower latitudes
(higher absolute latitude in the Southern Hemisphere convention) and the
profile becomes more symmetric.

\subsection{Magnitude of the shift}

For a body of half-width $a$ at depth $z$, the N--S displacement of the
TMI peak from the source centre is approximately
\begin{equation}
\Delta y \approx z \cot|I|.
\end{equation}
At Tower Hill, with $|I| = 70^\circ$ and a buried top at perhaps 200--400\,m
depth, the predicted shift is $\Delta y \approx 0.4\,a$ to $0.7\,a$,
roughly consistent with the observed displacement of a few hundred metres
to one kilometre visible in the maps.

% ─────────────────────────────────────────────────────────────────────────────
\section{Limitations and Caveats}
\label{sec:limitations}

\subsection{Remanent magnetisation}

The standard RTP assumes induced magnetisation parallel to the present-day
field.
Remanently magnetised rocks carry a frozen-in direction from the time of
cooling.
If the remanent direction differs significantly from the current inducing
field (which is common in Quaternary volcanics that have experienced
geomagnetic reversals or excursions), the standard RTP will:
\begin{itemize}
  \item fail to centre the anomaly correctly,
  \item potentially introduce phase errors larger than those it removes,
  \item produce an asymmetric ``corrected'' anomaly that looks worse than the
        original.
\end{itemize}
Where rock-magnetic data exist (e.g.\ from oriented drill-core samples),
the magnetisation direction can be incorporated explicitly.
The Tower Hill volcanics are young enough that their remanence direction
is unlikely to differ greatly from the current field, but this cannot be
assumed in general.

\subsection{Heterogeneous magnetisation}

The diatreme fill is a breccia: a chaotic mixture of country rock, basalt
clasts, and matrix.
The effective susceptibility and remanence vary laterally and vertically
in ways that a single RTP filter cannot capture.
The filter produces a spatially averaged correction that is valid only in
a bulk sense.

\subsection{Depth and susceptibility variations}

The forward model used in the notebook is a thin-sheet approximation at
a fixed depth.
Real diatremes have irregular geometry and variable fill thickness.
The RTP filter does not explicitly account for depth variation;
where source depth changes significantly across the map, the quality of the
phase correction will vary.

% ─────────────────────────────────────────────────────────────────────────────
\section{Connection to the Fourier Domain View}

The spectral decomposition in Part~B of the notebook directly motivates
the Tower Hill application:

\begin{enumerate}
  \item \textbf{Amplitude spectrum.}
        The log-amplitude spectrum of the Tower Hill TMI shows a roughly
        isotropic ring of elevated power at wavenumbers corresponding to
        wavelengths of 3--5\,km, consistent with the lateral scale of the
        diatreme.
        Higher wavenumbers carry noise and short-scale surface variation;
        lower wavenumbers are dominated by regional trends.

  \item \textbf{Phase spectrum.}
        The phase of the raw TMI spectrum is strongly asymmetric in the
        N--S direction, encoding the southward displacement of the positive
        lobe.
        After RTP, the phase structure becomes more isotropic at the
        wavenumbers carrying the maar signal.

  \item \textbf{RTP filter phase.}
        The correction $\angle H_\text{RTP}$ is direction-dependent:
        it is largest for N--S wavenumber components and nearly zero
        for purely E--W components.
        This anisotropy is a direct consequence of the N--S component
        of the horizontal field direction (the $\cos I\cos D$ term in
        $\Upsilon$) being larger than the E--W component at this declination.
\end{enumerate}

% ─────────────────────────────────────────────────────────────────────────────
\section{Key Takeaways}

\begin{itemize}
  \item A maar diatreme produces a positive TMI anomaly displaced northward
        from the source in the Southern Hemisphere due to the southward dip
        of the magnetic field.
  \item RTP is a Fourier-domain phase correction: it restores approximate
        symmetry and centres anomaly maxima over their sources without
        changing the amplitude spectrum.
  \item At $I \approx -70^\circ$, RTP is stable and reliable for Tower Hill;
        the filter becomes singular near the equator.
  \item The phase spectrum of the raw TMI encodes the geometric asymmetry;
        the RTP filter phase encodes the correction.
  \item Remanent magnetisation, heterogeneous susceptibility, and
        irregular geometry all limit the accuracy of RTP in practice.
\end{itemize}

\end{document}
